\section{Exact initial data and a moving-native reference}\label{sec:reference}
\subsection{Positive scalar majorant}\label{sec:majorant}
The exact initial deviations, with no first-order phase approximation, are
\begin{equation}\label{eq:initial}% RR17
\begin{gathered}
 u_0=(0,\cos\phi-1,0),\qquad b_0=(0,\sin\phi,0),\\
 \norm{Y(0)}_*^2
 =2\alpha[(\cos\phi-1)^2+\sin^2\phi]=4C .
\end{gathered}
\end{equation}
Write $y_0=2\sqrt C>0$. The known global strong solution and
\eqref{eq:forced}--\eqref{eq:nonlinearbound} give
$\norm{Y(t)}_*\le y_0+L\int_0^t\norm{Y(s)}_*^2\dd s=:Q(t)$.
On every finite interval, $Q$ is positive and continuously differentiable, with
$Q(0)=y_0$ and $Q'\le LQ^2$. Therefore
$(1/Q)'=-Q'/Q^2\ge-L$, and $1/Q(t)\ge1/y_0-Lt$.
Only while this lower bound is positive may it be inverted:
\begin{equation}\label{eq:majorant}% RR18
 \norm{Y(t)}_*\le Q(t)\le\frac{y_0}{1-Ly_0t},
 \qquad 0\le t<(Ly_0)^{-1}.
\end{equation}
Section~\ref{sec:constants} proves directly that $Ly_0t_F<1/25$ before this estimate is used up to $t_F$. There is no assumed bootstrap or new existence construction.

\subsection{Full reference, including static components}\label{sec:embedding}
Use $\dot Y^{(0)}=\mathcal A_0Y^{(0)}$ with precisely the initial data~\eqref{eq:initial}. Only the reference removes propagation and the quadratic remainder; actual $c$ and $\kappa$ stay positive.

Set $n=\ell/|\ell|$, $h_0=e^Tb_0=-\sin\phi/\sqrt2$, and
$b_\perp=b_0-h_0e=(\sin\phi/2,\sin\phi/2,0)$.
Then $Db_\perp=0$. The complete reference is
\begin{equation}\label{eq:embedding}% RR19
\begin{gathered}
 u^{(0)}(t)=u_0,\qquad b^{(0)}(t)=b_\perp+eX(t)/\sqrt{2\alpha},\\
 p^{(0)}(t)=\B n\,P(t)/\sqrt m,\quad
 r_c^{(0)}(t)=0,\quad z^{(0)}(t)=nZ(t)/\sqrt\beta .
\end{gathered}
\end{equation}
The real deviation, including its component parallel to $e$, and the inactive imaginary component are retained. Because $D^Tn=|\ell|e$ and
$y^{(0)}=\sqrt m\,nP$, all blocks of $\mathcal A_0$ reduce to
\begin{equation}\label{eq:oscillator}% RR20
\begin{gathered}
 \dot X=-\omega_NP,\qquad
 \dot P=\omega_NX-\omega_RZ,\qquad \dot Z=\omega_RP,\\
 X_0=-\sqrt\alpha\sin\phi<0,\qquad P_0=Z_0=0 .
\end{gathered}
\end{equation}
The negative sign follows from the exact phase launch, with $0<\phi<\pi/2$.
The solution is
\begin{equation}\label{eq:solution}% RR21
\begin{aligned}
 X(t)&=X_0\frac{\omega_R^2+\omega_N^2\cos(\varpi t)}{\varpi^2},\\
 P(t)&=X_0\frac{\omega_N}{\varpi}\sin(\varpi t),\\
 Z(t)&=X_0\frac{\omega_N\omega_R}{\varpi^2}
                   [1-\cos(\varpi t)] .
\end{aligned}
\end{equation}
Differentiation verifies the three ODEs and all initial data, and
$X^2+P^2+Z^2=X_0^2$. Consequently,
\begin{equation}\label{eq:referencezeros}% RR22
\begin{gathered}
 (p^{(0)},r_c^{(0)})(t_R)=0,\qquad
 (p^{(0)},r_c^{(0)},z^{(0)})(t_F)=0,\\
 \frac{X_0^2}{2}=C\cos^2(\phi/2),\qquad
 \frac{|Z(t_R)|}{\sqrt{2C}}
 =\frac{2\sqrt{30}}{17}\cos(\phi/2).
\end{gathered}
\end{equation}
In particular, $|X_0|/\sqrt{2C}=\cos(\phi/2)\le1$, with positive cosine. This reference is a strong curve: a fixed $H^1$ profile multiplies a smooth scalar. Its period is not asserted to be an actual nonlinear period, and the actual trajectory is not restarted at either landmark.
